  % Over a century ago, \citet{1933AcHPh...6..110Z} postulated the existence of dark matter to explain the excessively large velocity dispersions of galaxies within their galaxy clusters. Since then, mounting evidence has been assembled to support this claim. \citet{1970IAUS...38..351F,1970ApJ...161..802F} elaborated on the dynamics of disk galaxies and briefly acknowledged in the appendix of his paper that disk models cannot fully explain the rotation curves of NGC~300 and M~33 and that there must be the same amount of mass or more that is undetected and beyond the extend of the galaxies to explain their kinematics. Concurrently, \citet{1970ApJ...159..379R} studied the rotation of Andromeda and observed it to be flat---not decreasing as one would expect from a gravitational potential with more mass concentration in the center. Eventually, astronomers returned to Zwicky's non-luminous matter to explain these fast-rotating disks, as discussed in \citet{1983SciAm.248f..96R}'s review.
  
  % Beyond the Galactic scale, dark matter is essential for explaining many cosmological phenomena. For instance, the Planck mission \citep{2020A&A...641A...1P} was designed to study the Cosmic Microwave Background (CMB), whose spectrum closely follows that of a black body with a peak temperature in the millimeter wavelength range, corresponding to approximately 2.725~K. However, there are slight temperature variations on the order of one part in ten thousand. Notably, the colder regions of the CMB temperature map correlate with whole-sky gravitational lensing maps \citep{2020A&A...641A...8P}--since the photons lost energy as they escaped the gravitational wells. The $\Lambda$~Cold Dark Matter ($\Lambda$CDM) model accounts for these observations almost perfectly, proposing that dark matter exists and that its distribution originated from stochastic quantum fluctuations within the first second of the universe, which led to density variations. These variations propagated as the universe expanded, shaping the large-scale structure and distribution of galaxies we observe today.

  % As laid out in \citet{2021PrPNP.11903865A}'s review, despite significant progress, the true nature of dark matter remains elusive, and scientists continue to search for evidence of its existence beyond its gravitational effects on cosmic and Galactic scales. A key focus today is detecting dark matter on sub-galactic scales because, at these scales, different cosmological scenarios predict different lumpiness of the halos.
    
  % \citet{2019Galax...7...81Z} provides a review of the different proposed mechanisms of forming dark matter haloes theoretically and how to model them practically, showcasing how different assumptions for the fundamental particle behavior of dark matter propagate to the mass and size distribution of sub-haloes on the galactic scale. The authors rightly state that until particles are discovered, $\Lambda$CDM remains a hypothesis. Since the true nature of dark matter is elusive, the possibilities for sub-galactic behavior and structure of dark matter are vast. In turn, obtaining constraints on the distribution of dark matter subhaloes will constrain the particle nature of dark matter. As proposed by \citet{2012ApJ...748...20C}, stellar streams are the most promising tool to limit the number and mass density of the sub-haloes. 

    % \citet{2022MNRAS.516.5331M} showed that many metal-poor streams exist in the Milky Way. On the other hand, it is possible that some globular clusters do not lose a significant number of stars to create tidal tails. For instance, \citet{2022A&A...667A.112V} demonstrated that clusters existing in dark matter subhaloes could place the tidal radius far beyond their half-mass radii and shield stars from escape. Another factor is indeed detectability. Additionally, there are many clusters in the internal region of the galaxy, a region over which \texttt{streamfinder} was not applied since the high stellar density implies a computational cost that is far too expensive. Also, \citet{2018MNRAS.474.2479B} demonstrates that internal dynamics preferentially kick out small stars first, which have lower luminosities. Earlier escape times place larger distances to the cluster, making detection more difficult. 

        %Our results are consistent with the findings from \citet{2021MNRAS.505.5978V}, who identified a wide range of orbital solutions for the least bound clusters, including retrograde and prograde orbits. In our potential model, only three cluster samplings had positive orbital energies.

        %: Slidr, Sylgr, Ylgr, Fimbulthul, Svol, Gjoll, Leiptr, Phlegeton and Fjorm.  According to the most recent estimates provided by \citet{2020MNRAS.493.4978S}, Fjorm, in particular, has an extension of more than 150 degrees on the sky, which has a projected length more than five times greater than that of Palomar~5. 

            % Each group can be simplified as one-dimensional collisionless dispersion. These packets of stars can be described as a population with the same initial position that drifts about a mean velocity, $\langle v \rangle$, with velocity dispersion, $\sigma_v$. The density evolution is a Gaussian whose amplitude decreases inversely with time, and whose standard deviation grows linearly with time:
    % \begin{equation}    
          % \rho(x,t) = \frac{1}{\sigma_v t \sqrt{2\pi} }\textrm{exp}\left(-\frac{1}{2}\left(\frac{x-\langle v \rangle t}{\sigma_v t}\right)^2\right).
    % \end{equation}

    % It is important to note that within a single packet, the stars self-segregate. The position of the density peak is at $\langle v \rangle t$, while the faster stars are ahead and the slower ones are behind. The packets are separated by a distance of $\langle  v\rangle T$, where $T$ is the orbital period. Regarding our simulations, since the mass and radius of the host are held constant, the mean velocity and velocity dispersion are the same across all packets. Secondly, the density of each packet decreases since the easier-to-remove stars are stripped early on. 

    % Consider an impact that happens at $x\prime$. The distance from the impact site to each packet is: $\delta x= x\prime - \langle v\rangle \left(t-n T\right)$, where $n$ is the packet index, with $n=0$ being the first packet escaped at the first pericenter passage in the simulation. Note that $\delta x\propto n$, and velocity segregation means $v\propto x$. Therefore, the different parts of the separate packets are hit at a given impact site. The earlier packets are hit in regions where the stars have lower velocities than the later packets. As a result, the gaps in the later packets quickly overtake those from the earlier packets. As a result, the gaps across the packets go out of phase and interfere destructively.

	  % The aforementioned analysis highlights the complexity of the conditions on the stream that are necessary for allowing gap formation. Despite this, in the attempts to gain further preliminary insight, we can create an even simpler consideration of a gap created in two different pockets, considering the gaps' drift rates and width, ignoring their density contrasts. If we consider two packets with the same dimensions $\sigma_x$ and an offset of $\delta x$, the gaps are out of phase when the separation is greater than the width: $1\sim<\sigma_x/\delta x$. If $\delta x$ grows much faster than the gap can develop, then the gap is never visible. If the gaps grow in size much faster than their drift rate, then they persist and continue to grow in significance. In the simulations presented in this analysis, we only observe the two extremes and do not produce a case where a significant gap is created to blur out. However, this is not necessarily the general case.

    % Nonetheless, this is a first look at considering the distribution of stream stars originating from shocked outflow, and they allow or disallow gap formation. Fuller considerations would include mass loss from internal dynamics, such as three-body encounters, and a non-constant mass loss rate due to changes in the density distribution resulting from internal mass evolution.
